\section*{Capítulo 4: Cómo doblar el espaciotiempo}
\begin{description}[style=multiline,leftmargin=2em,labelwidth=1.5em]

% DOCX paragraph 2032
\item[1.]\hypertarget{note:04-01}{}Como recordarás del capítulo 2, la constante
gravitatoria \emph{G} es igual a
6,67408\(\times\)10\textsuperscript{−11} m\textsuperscript{3}/(kg
s\textsuperscript{2}), donde m, kg y s representan, respectivamente, metros,
kilogramos y segundos, que son unidades del Sistema Internacional de Unidades
(SI). De ello se deduce que la cantidad \emph{GM/(c}\textsuperscript{2} \emph{R}) no
tiene unidades, pues todas se cancelan. Por tanto, la cantidad
\emph{GM/(c}\textsuperscript{2} \emph{R}) es un \emph{número puro}.

% DOCX paragraph 2033
\item[2.]\hypertarget{note:04-02}{}El símbolo \(\oplus\) se utiliza en astronomía
para indicar el planeta Tierra; así, por ejemplo, \(M_{\oplus}\) indica la masa de
nuestro planeta. Del mismo modo, el símbolo \(\odot\), que utilizaremos enseguida,
indica el Sol y cualquier cantidad relativa a él.

% DOCX paragraph 2034
\item[3.]\hypertarget{note:04-03}{}El pársec es una unidad de longitud que se utiliza
habitualmente en astronomía, pues permite expresar las enormes distancias que
encontramos fuera del sistema solar. Por este motivo, en astronomía nunca se usan
kilómetros, sino años luz y sus múltiplos. Recuerdo que un «año luz» no mide un
intervalo de tiempo, sino de espacio: el espacio que recorre un haz de luz durante
un año. Expresado en kilómetros, corresponde a diez billones
(10\textsuperscript{13}) de kilómetros. Por su parte, un pársec corresponde a
3,26 años luz, es decir, unos treinta billones de kilómetros. Está claro que es
mucho más fácil hablar de un pársec que de decenas de billones de kilómetros…

\end{description}
